%
% Copyright (C) 1995, 2000, 2013 American Mathematical Society.
% Copyright (C) 2016-2025 LaTeX Project and American Mathematical Society.
%
% This work may be distributed and/or modified under the
% conditions of the LaTeX Project Public License, either version 1.3c
% of this license or (at your option) any later version.
% The latest version of this license is in
%   https://www.latex-project.org/lppl.txt
% and version 1.3c or later is part of all distributions of LaTeX
% version 2005/12/01 or later.
%
% This work has the LPPL maintenance status `maintained'.
%
% The Current Maintainer of this work is the LaTeX Project.
%
\documentclass[leqno,titlepage,openany]{amsldoc}[1999/12/13]

\def\MaintainedByLaTeXTeam#1{%
\begin{center}%
\fbox{\fbox{\begin{tabular}{@{}l@{}}%
This file is maintained by the \LaTeX{} Project team.\\%
Bug reports can be opened (category \texttt{#1}) at\\%
\url{https://latex-project.org/bugs/}.\end{tabular}}}\end{center}}

\ifx\UndEfiNed\url
  \ClassError{amsldoc}{%
    This version of amsldoc.tex must be processed\MessageBreak
    with a newer version of amsldoc.cls (2.02 or later).}{}
\fi

\renewcommand\bibname{References}

%    Definitions that should be moved to amsldoc.cls
%    when it is next updated
\makeatletter
\let\cleardouble@page\cleardoublepage
\AtBeginDocument{%
  \ifx\cleardouble@page\cleardoublepage
    \def\cleardoublepage{\clearpage{\pagestyle{empty}\cleardouble@page}}
  \fi
}

% #1 is empty or one of class, environment, option, package
% could use different font in each case
\makeatletter
\def\@category@index#1#2{%
    \autoindex{%
        #2\actualchar\string\index#1font{#2}%
        \ifx\@nil#1\@nil\else\space#1\fi
    }%
}

\let\indexpackagefont\textsf
\let\indexclassfont\textsf
\let\indexoptionfont\texttt
\let\indexenvironmentfont\texttt
\let\indexfont\texttt
\makeatother
%    For a package that won't be indexed, i.e. amsmath
\newcommand{\nipkg}{\textsf}
\def\tikz/{Ti\textit{k}Z}
%    End definitions for amsldoc.cls

\title{User's Guide for the \nipkg{amsmath} Package (Version~2.1)}
\author{American Mathematical Society, \LaTeX\ Project}
\date{1999-12-13\\(revised 2002-02-25, 2016-11-14, 2018-04-05, 2019-10-14, 2020-02-18)}
\makeatletter
\def\@thanks{\bigskip\MaintainedByLaTeXTeam{amslatex}}
\makeatother

% Very minor changes are not listed in the \date above:
%
% 2021-08-22: reported typo fixed mathools -> mathtools
% 2025-05-23: drop spurious option list


%    Use the amsmath package and amscd package in order to print
%    examples.
\usepackage{amsmath}
\usepackage{amscd}

\usepackage{imakeidx}
%    Need alternate coding for vertical bar to get it into index.
%    Based on \qcbang and \cnbang in amsldoc.cls, refined by David Carlisle.
%    Give it a sort field that will be at the end of the symbols;
%    must avoid " and | which have special meanings.
\DeclareRobustCommand{\qcvert}{\qc\|\index{{z1}@{\ntt\char'174}}}
\DeclareRobustCommand{\cnvert}{%
  \ncn{\|}\index{{z0}@{\ntt\bslchar\char'174}}}
%    Need alternate index sort field for bang to allow & to sort first
\DeclareRobustCommand{\altcnbang}{%
  \ncn{\!}\index{"!b@{\ntt\bslchar\qcbang}}}
\DeclareRobustCommand{\cnamp}{%
  \index{"!a@\texttt{\&}}}

\makeindex[intoc] % generate index data
\providecommand{\see}[2]{\textit{see} #1}

%    The amsldoc class includes a number of features useful for
%    documentation about TeX, including:
%
%    ---Commands \tex/, \amstex/, \latex/, ... for uniform treatment
%    of the various logos and easy handling of following spaces.
%
%    ---Commands for printing various common elements: \cn for command
%    names, \fn for file names (including font-file names), \env for
%    environments, \pkg and \cls for packages and classes, etc.

%    Many of the command names used here are rather long and will
%    contribute to poor linebreaking if we follow the \latex/ practice
%    of not hyphenating anything set in tt font; instead we selectively
%    allow some hyphenation.
\allowtthyphens % defined in amsldoc.cls

\hyphenation{ac-cent-ed-sym-bol add-to-counter add-to-length align-at
  aligned-at allow-dis-play-breaks ams-art ams-cd ams-la-tex amsl-doc
  ams-symb ams-tex ams-text ams-xtra bmatrix bold-sym-bol cen-ter-tags
  eqn-ar-ray idots-int int-lim-its latex med-space neg-med-space
  neg-thick-space neg-thin-space no-int-lim-its no-name-lim-its
  over-left-arrow over-left-right-arrow over-right-arrow pmatrix
  qed-sym-bol set-length side-set small-er tbinom the-equa-tion
  thick-space thin-space un-der-left-arrow un-der-left-right-arrow
  un-der-right-arrow use-pack-age var-inj-lim var-proj-lim vmatrix
  xalign-at xx-align-at}

%    Prepare for illustrating the \vec example
\newcommand{\vect}[1]{\mathbf{#1}}

\newcommand{\booktitle}[1]{\textit{#1}}
\newcommand{\journalname}[1]{\textit{#1}}
\newcommand{\seriesname}[1]{\textit{#1}}

%    Command to insert and index a particular phrase. Doesn't work for
%    certain kinds of special characters in the argument.
\newcommand{\ii}[1]{#1\index{#1}}

\newcommand{\vstrut}[1]{\vrule width0pt height#1\relax}

%    An environment for presenting comprehensive address information:
\newenvironment{infoaddress}{%
  \par\topsep\medskipamount
  \trivlist\centering
  \item[]%
  \begin{minipage}{.7\columnwidth}%
  \raggedright
}{%
  \end{minipage}%
  \endtrivlist
}

\newenvironment{eqxample}{%
  \par\addvspace\medskipamount
  \noindent\begin{minipage}{.5\columnwidth}%
  \def\producing{\end{minipage}\begin{minipage}{.5\columnwidth}%
    \hbox\bgroup\kern-.2pt\vrule width.2pt%
      \vbox\bgroup\parindent0pt\relax
%    The 3pt is to cancel the -\lineskip from \displ@y
    \abovedisplayskip3pt \abovedisplayshortskip\abovedisplayskip
    \belowdisplayskip0pt \belowdisplayshortskip\belowdisplayskip
    \noindent}
}{%
  \par
%    Ensure that a lonely \[\] structure doesn't take up width less than
%    \hsize.
  \hrule height0pt width\hsize
  \egroup\vrule width.2pt\kern-.2pt\egroup
  \end{minipage}%
  \par\addvspace\medskipamount
}

%    Errors in output shouldn't be \texttt, and they may also be longer
%    than one line, but they should be identified by the "errorbullet".
%    (They were formerly labeled as aubsections.)
\newenvironment{erroro}[1]{%
  \par\addvspace\medskipamount
  \noindent\hangindent\parindent
  \hbox to\parindent{\errorbullet\hfil}\ignorespaces
  #1\par\vspace{\smallskipamount}
}{%
  \par\addvspace\medskipamount
}

%    The chapters are so short, perhaps we shouldn't call them by the
%    name `Chapter'. We make \chaptername read an argument in order to
%    remove a following \space or "{} " (both possibilities are present
%    in book.cls).

\renewcommand{\chaptername}[1]{}
\newcommand{\chapnum}[1]{\mdash #1\mdash }
\makeatletter
\def\@makechapterhead#1{%
  \vspace{1.5\baselineskip}%
  {\parindent \z@ \raggedright \reset@font
    \ifnum \c@secnumdepth >\m@ne
      \large\bfseries \chapnum{\thechapter}%
      \par\nobreak
      \vskip.5\baselineskip\relax
    \fi
    #1\par\nobreak
    \vskip\baselineskip
  }}
\makeatother

%    A command for ragged-right parbox in a tabular.
\newcommand{\rp}{\let\PBS\\\raggedright\let\\\PBS}

%    Non-indexed file name
\newcommand{\nfn}[1]{\texttt{#1}}

%    For the examples in the math spacing table.
%%\newcommand{\lspx}{\mbox{\rule{5pt}{.6pt}\rule{.6pt}{6pt}}}
%%\newcommand{\rspx}{\mbox{\rule[-1pt]{.6pt}{7pt}%
%%  \rule[-1pt]{5pt}{.6pt}}}
\newcommand{\lspx}{\mathord{\Rightarrow\mkern-1mu}}
\newcommand{\rspx}{\mathord{\mkern-1mu\Leftarrow}}
\newcommand{\spx}[1]{$\lspx #1\rspx$}

%    For a list of characters representing document input.
\newcommand{\clist}[1]{%
  \mbox{\ntt\spaceskip.2em plus.1em \xspaceskip\spaceskip#1}}

%    Fix weird \latex/ definition of rightmark.
\makeatletter
\def\rightmark{\expandafter\@rightmark\botmark{}{}}
%    Also turn off section marks.
\let\sectionmark\@gobble
\renewcommand{\chaptermark}[1]{%
  \uppercase{\markboth{\rhcn#1}{\rhcn#1}}}
\newcommand{\rhcn}{\thechapter. }
\makeatother

%    Include down to \section but not \subsection, in toc:
\setcounter{tocdepth}{1}

%    Subheadings for the bibliography
\newif\iffirstbibsubhead \firstbibsubheadtrue
\newcommand{\bibsubhead}[1]{%
  \iffirstbibsubhead \firstbibsubheadfalse
  \else \addvspace{\medskipamount}
  \fi
  \item[]\hspace*{-\leftmargin}\textbf{#1}\par
  \vspace{\smallskipamount}
}

\DeclareMathOperator{\ix}{ix}
\DeclareMathOperator{\nul}{nul}
\DeclareMathOperator{\End}{End}
\DeclareMathOperator{\xxx}{xxx}

\usepackage{hyperref}
\hypersetup{pdfborder=0 0 0}
\makeatletter
\let\oldcs\cs
\def\cs#1{\texorpdfstring{\oldcs{#1}}{\@backslashchar\@backslashchar#1}}
\let\cn\cs
\makeatother
\begin{document}

%%%%%%%%%%%%%%%%%%%%%%%%%%%%%%%%%%%%%%%%%%%%%%%%%%%%%%%%%%%%%%%%%%%%%%%%
\frontmatter

\maketitle

\pagestyle{headings}
\tableofcontents
\cleardoublepage % for better page number placement

%%%%%%%%%%%%%%%%%%%%%%%%%%%%%%%%%%%%%%%%%%%%%%%%%%%%%%%%%%%%%%%%%%%%%%%%
\mainmatter
\chapter{Introduction}

\index{amsmath@\texttt{amsmath} package|(}

The \nipkg{amsmath} package is a \LaTeX{} package that provides
miscellaneous enhancements for improving the information structure and
printed output of documents that contain mathematical formulas. Readers
unfamiliar with \LaTeX{} should refer to \cite{lamport}. If you have an
up-to-date version of \LaTeX{}, the \nipkg{amsmath} package is normally
provided along with it. Upgrading when a newer version of the
\nipkg{amsmath} package is released can be done via
\url{http://mirror.ctan.org/macros/latex/required/amsmath.zip}.

This documentation describes the features of the \nipkg{amsmath} package
and discusses how they are intended to be used. It also covers some
ancillary packages:
\begin{ctab}{lll}
\pkg{amsbsy}& \pkg{amsopn}& \pkg{amsxtra}\\
\pkg{amscd}& \pkg{amstext}\\
\end{ctab}
These all have something to do with the contents of math formulas. For
information on extra math symbols and math fonts, see \cite{amsfonts}
and \url{https://www.ams.org/tex/amsfonts.html}. For documentation of the
\pkg{amsthm} package or AMS document classes (\cls{amsart},
\cls{amsbook}, etc.\@) see \cite{amsthdoc} or \cite{amshandbk} and
\url{https://www.ams.org/tex/author-info.html}.

If you are a long-time \latex/ user and have lots of mathematics in what
you write, then you may recognize solutions for some familiar problems
in this list of \nipkg{amsmath} features:
\begin{itemize}

\item A convenient way to define new `operator name' commands analogous
to \cn{sin} and \cn{lim}, including proper side spacing and automatic
selection of the correct font style and size (even when used in
sub- or superscripts).

\item Multiple substitutes for the \env{eqnarray} environment to make
various kinds of equation arrangements easier to write.

\item Equation numbers automatically adjust up or down to avoid
overprinting on the equation contents (unlike \env{eqnarray}).

\item Spacing around equals signs matches the normal spacing in the
\env{equation} environment (unlike \env{eqnarray}).

\item A way to produce multiline subscripts as are often used with
summation or product symbols.

\item An easy way to substitute a variant equation number for a given
equation instead of the automatically supplied number.

\item An easy way to produce subordinate equation numbers of the form
(1.3a) (1.3b) (1.3c) for selected groups of equations.

\end{itemize}

The \nipkg{amsmath} package is distributed together with some small
auxiliary packages:
\begin{description}
\item[\nipkg{amsmath}] Primary package, provides various features for
  displayed equations and other mathematical constructs.

\item[\pkg{amstext}] Provides a \cn{text} command for
  typesetting a fragment of text inside a display.

\item[\pkg{amsopn}] Provides \cn{DeclareMathOperator} for defining new
  `operator names' like \cn{sin} and \cn{lim}.

\item[\pkg{amsbsy}] For backward compatibility this package continues
to exist but use of the newer \pkg{bm} package that comes with \LaTeX{}
is recommended instead.

\item[\pkg{amscd}] Provides a \env{CD} environment for simple
  commutative diagrams (no support for diagonal arrows).

\item[\pkg{amsxtra}] Provides certain odds and ends such as
  \cn{fracwithdelims} and \cn{accentedsymbol}, for compatibility with
  documents created using version 1.1.

\end{description}

The \nipkg{amsmath} package incorporates \pkg{amstext}, \pkg{amsopn}, and
\pkg{amsbsy}. The features of \pkg{amscd} and \pkg{amsxtra}, however,
are available only by invoking those packages separately.

The independent \pkg{mathtools} package \cite{mt} provides some
enhancements to \nipkg{amsmath}; \pkg{mathtools} loads \nipkg{amsmath}
automatically, so there is no need to separately load \nipkg{amsmath}
if \pkg{mathtools} is used.  Some \pkg{mathtools} facilities
will be noted below as appropriate.

%%%%%%%%%%%%%%%%%%%%%%%%%%%%%%%%%%%%%%%%%%%%%%%%%%%%%%%%%%%%%%%%%%%%%%%%

\chapter{Options for the \nipkg{amsmath} package}\label{options}

The \nipkg{amsmath} package has the following options:
\index{options!amsmath@\texttt{amsmath} package options|(}
\begin{description}

\item[\opt{centertags}] (default) For an equation containing a
\env{split} environment, place equation
numbers\index{equationn@equation numbers!vertical placement}
vertically centered on the total height of the equation.

\item[\opt{tbtags}] `Top-or-bottom tags': For an equation containing
a \env{split} environment, place equation
numbers\index{equationn@equation numbers!vertical placement} level with the last
(resp.\ first) line, if numbers are on the right (resp.\ left).

\item[\opt{sumlimits}] (default) Place the subscripts and
superscripts\index{subscripts and superscripts!placement}\relax
\index{limits|see{subscripts and superscripts}} of summation symbols
above and below, in displayed equations. This option also affects other
symbols of the same type\mdash $\prod$, $\coprod$, $\bigotimes$,
$\bigoplus$, and so forth\mdash but excluding integrals (see below).

\item[\opt{nosumlimits}] Always place the subscripts and superscripts of
summation-type symbols to the side, even in displayed equations.

\item[\opt{intlimits}] Like \opt{sumlimits}, but for
integral\index{integrals!placement of limits} symbols.

\item[\opt{nointlimits}] (default) Opposite of \opt{intlimits}.

\item[\opt{namelimits}] (default) Like \opt{sumlimits}, but for certain
`operator names' such as $\det$, $\inf$, $\lim$, $\max$, $\min$, that
traditionally have subscripts%
\index{subscripts and superscripts!placement} placed underneath when
they occur in a displayed equation.

\item[\opt{nonamelimits}] Opposite of \opt{namelimits}.

\end{description}

\goodbreak

To use one of these package options, put the option name in the optional
argument of the \cn{usepackage} command\mdash e.g.,
\verb"\usepackage[intlimits]{amsmath}".
For AMS document classes and any other classes that preload \nipkg{amsmath}
desired options must be specified with the \cn{documentclass}\mdash
e.g.,\\
\verb+\documentclass[intlimits,tbtags,reqno]{amsart}+.

The \nipkg{amsmath} package also recognizes the following options which
are normally selected (implicitly or explicitly) through the
\cn{documentclass} command, and thus need not be repeated in the option
list of the \cn{usepackage}|{amsmath}| statement.
\begin{description}

\item[\opt{leqno}] Place equation numbers on the left.%
\index{equationn@equation numbers!left or right placement}

\item[\opt{reqno}] Place equation numbers on the right.

\item[\opt{fleqn}] Position equations at a fixed indent from the left
margin rather than centered in the text column.%
\index{displayed equations!centering}

\end{description}

Three options have been added to control the space to the left of
\env{aligned} and \env{gathered} environments. Prior to the 2017 release
a thin space was added to the left but not the right of these constructs.
This appears to have been an accidental feature of the implementation and
was typically corrected by prefixing the environments with \altcnbang.

The new default behavior is aimed to ensure that the environments do
not have a thin space added in most cases, and that existing documents
using |\!\begin{aligned}| continue to work as before.

\begin{description}
\item[\opt{alignedleftspaceyes}] Always add \cn{\,} to the left of \env{aligned} and \env{gathered}.
\item[\opt{alignedleftspaceno}] Never add \cn{\,} to the left of \env{aligned} and \env{gathered}.
\item[\opt{alignedleftspaceyesifneg}] Only add \cn{\,} if the environment is prefixed by negative space. (New default behavior.)
\end{description}

\index{options!amsmath@\texttt{amsmath} package options|)}


%%%%%%%%%%%%%%%%%%%%%%%%%%%%%%%%%%%%%%%%%%%%%%%%%%%%%%%%%%%%%%%%%%%%%%%%
\end{document}
